\section{Ablation Study}
\label{sec:ablation}

\subsection{Overview}

The primary controlled ablation in this assignment is the comparison between \textbf{EXP-2 (frozen backbone)} and \textbf{EXP-3 (full fine-tune)}. These two experiments share identical architectures, datasets, splits, batch sizes, optimizer type, and loss function. The \textbf{only changed factor} is the set of trainable parameters: in EXP-2, only the 594,437-parameter classification head is updated; in EXP-3, all 66,957,317 parameters are updated.

This design directly measures the contribution of backbone adaptation to downstream sentiment accuracy.

\subsection{Results: Frozen vs.\ Full Fine-Tuning}

\begin{table}[H]
    \centering
    \caption{Ablation: effect of freezing vs.\ training the DistilBERT backbone.}
    \label{tab:ablation_frozen}
    \begin{tabularx}{\textwidth}{l r r r r r r}
        \toprule
        \textbf{Setting} & \textbf{Trainable Params} & \textbf{Train Time (s)} & \textbf{Test Acc.} & \textbf{Macro-F1} \\
        \midrule
        EXP-2: Frozen backbone (head only) & 594,437   & 561.29   & 38.94\% & 0.3779 \\
        EXP-3: Full fine-tune              & 66,957,317 & 1,019.09 & 85.33\% & 0.8577 \\
        \midrule
        \textbf{Difference (EXP-3 $-$ EXP-2)} & +66.4M & +457.80 & \textbf{+46.39 pts} & \textbf{+0.4798} \\
        \bottomrule
    \end{tabularx}
\end{table}

The ablation reveals a dramatic 46.39-point accuracy gain from unfreezing the backbone. Full fine-tuning takes 1.82$\times$ longer to train but yields a 2.27$\times$ higher macro-F1.

\subsection{Efficiency Gains from Frozen Training}

Despite its lower accuracy, the frozen-backbone configuration offers two practical advantages:
\begin{itemize}
    \item \textbf{Training speed:} 331 seconds faster (561 vs.\ 1,019 seconds) due to no gradient computation through the large backbone.
    \item \textbf{GPU memory:} Significantly lower memory footprint during training, as no gradient tensors are stored for the 66.4M frozen parameters. This enables larger batch sizes or running on GPUs with limited VRAM.
\end{itemize}

\subsection{Caveat: Confounded Learning Rate}

The frozen-backbone experiment (EXP-2) used a learning rate of $5\times10^{-3}$, which was chosen for BiLSTM-scale training. In the full fine-tune setting (EXP-3), a much smaller learning rate of $2\times10^{-5}$ was used, appropriate for adapting a large pretrained model.

For head-only training, the optimal learning rate is typically an order of magnitude larger than for full fine-tuning (since the backbone gradient signal does not need to be preserved), but the specific value of $5\times10^{-3}$ may still have been suboptimal. A systematic sweep over learning rates for EXP-2 (e.g.\ $\{1\times10^{-4},\,5\times10^{-4},\,1\times10^{-3}\}$) might yield meaningfully better frozen-backbone performance. As such, the EXP-2 vs.\ EXP-3 comparison is not a perfectly fair head-to-head ablation of ``frozen vs.\ full'' — it is confounded by the learning rate choice.

\subsection{EXP-1 is Not a Single-Factor Ablation}

We note that the EXP-1 BiLSTM cannot be treated as a clean single-factor ablation relative to EXP-2/3: it differs in architecture (recurrent vs.\ Transformer), pretraining (none vs.\ pretrained), learning rate, weight decay, batch size, and max epochs. Its role is to establish a practical baseline rather than to isolate a single variable.
